% Chapter 1: Exploratory Data Analysis
\chapter{Exploratory Data Analysis}

\section{Introduction}
This project investigates forest fires in Algeria and Tunisia through a data mining approach. We integrate five datasets---fire records, land cover, climate, elevation, and soil---to understand fire occurrence factors. The objectives are:
\begin{enumerate}
    \item Perform exploratory data analysis on environmental datasets
    \item Build predictive models for fire occurrence
    \item Identify spatial patterns and fire-prone clusters
\end{enumerate}

The repository structure and file descriptions are provided in Appendix~\ref{app:project}.

\section{Fire Dataset}
The fire dataset is sourced from NASA's FIRMS (Fire Information for Resource Management System), providing active fire detections from MODIS and VIIRS satellites.

\textbf{Key attributes:} Latitude/Longitude (location), Brightness (fire intensity in Kelvin), FRP (Fire Radiative Power in MW), Confidence level, and Acquisition date/time.

\begin{figure}[H]
    \centering
    \fbox{\parbox{0.85\textwidth}{\centering\vspace{2cm}\textbf{[PLACEHOLDER: Fire Dataset Statistics and Structure]}\vspace{2cm}}}
    \caption{Fire dataset structure and descriptive statistics}
    \label{fig:fire_structure}
\end{figure}

\textbf{Temporal patterns:} Fire activity peaks during summer months (June-September), with August recording the highest detections. This aligns with Mediterranean climate characteristics---hot, dry summers with minimal precipitation.

\textbf{Spatial distribution:} Fires concentrate in northern coastal regions with Mediterranean vegetation, mountainous forested areas (Atlas Mountains), and transition zones between Mediterranean and semi-arid climates.

\begin{figure}[H]
    \centering
    \fbox{\parbox{0.85\textwidth}{\centering\vspace{3cm}\textbf{[PLACEHOLDER: Fire Spatial Distribution Map]}\vspace{3cm}}}
    \caption{Geographic distribution of fire detections across Algeria and Tunisia}
    \label{fig:fire_spatial}
\end{figure}

\section{Land Cover Dataset}
Land cover data from the FAO Global Land Cover Share database classifies surface types including forests, shrublands, grasslands, croplands, urban areas, bare land, and water bodies.

\textbf{Key finding:} Forests and shrublands show the highest fire incidence due to fuel availability. Bare lands and urban areas have minimal fire activity. This relationship guided our decision to include land cover as a feature in modeling.

\begin{figure}[H]
    \centering
    \fbox{\parbox{0.85\textwidth}{\centering\vspace{3cm}\textbf{[PLACEHOLDER: Land Cover Map and Fire Distribution by Type]}\vspace{3cm}}}
    \caption{Land cover classification and fire occurrence by vegetation type}
    \label{fig:landcover}
\end{figure}

\section{Climate Dataset}
Climate data comes from CHIRPS (precipitation) and TerraClimate (temperature) with monthly resolution at approximately 4~km spatial resolution.

\textbf{Variables:} Precipitation (mm/month), Maximum temperature (°C), Minimum temperature (°C).

\textbf{Climate-fire relationship:} The Mediterranean climate of the study region shows strong correlation with fire patterns:
\begin{itemize}
    \item \textbf{Temperature:} Positive correlation with fire occurrence
    \item \textbf{Precipitation:} Negative correlation---drought conditions elevate risk
    \item \textbf{Decision:} We derived a drought index combining temperature and precipitation effects
\end{itemize}

\begin{figure}[H]
    \centering
    \fbox{\parbox{0.85\textwidth}{\centering\vspace{3cm}\textbf{[PLACEHOLDER: Climate Patterns and Fire Correlation]}\vspace{3cm}}}
    \caption{Seasonal climate patterns and correlation with fire occurrence}
    \label{fig:climate}
\end{figure}

\section{Elevation Dataset}
Elevation data from GMTED2010 provides a Digital Elevation Model at 15 arc-second resolution (approximately 500~m).

\textbf{Topographic analysis:}
\begin{itemize}
    \item \textbf{Slope:} Steeper slopes promote faster fire spread (preheating of upslope fuels)
    \item \textbf{Aspect:} South-facing slopes are drier due to solar radiation
    \item \textbf{Decision:} We derived terrain complexity and coastal proximity as additional features
\end{itemize}

\begin{figure}[H]
    \centering
    \fbox{\parbox{0.85\textwidth}{\centering\vspace{3cm}\textbf{[PLACEHOLDER: Elevation Map and Topographic Analysis]}\vspace{3cm}}}
    \caption{Digital elevation model and derived topographic features}
    \label{fig:elevation}
\end{figure}

\section{Soil Dataset}
Soil data from the Harmonized World Soil Database v2.0 (HWSD2) provides properties including organic carbon, nitrogen, pH, and texture fractions (sand, clay, silt).

\textbf{Soil-fire relationship:}
\begin{itemize}
    \item \textbf{Organic carbon:} Higher content supports denser vegetation, increasing fuel loads
    \item \textbf{Soil texture:} Sandy soils drain faster, leading to drier conditions
    \item \textbf{Decision:} We included soil properties as secondary predictors
\end{itemize}

\begin{figure}[H]
    \centering
    \fbox{\parbox{0.85\textwidth}{\centering\vspace{3cm}\textbf{[PLACEHOLDER: Soil Properties Distribution]}\vspace{3cm}}}
    \caption{Soil classification and properties at fire locations}
    \label{fig:soil}
\end{figure}

\section{Summary}
The exploratory analysis reveals:
\begin{itemize}
    \item \textbf{Fire Dataset:} Strong seasonal patterns (summer peak) and spatial concentration in northern forested regions
    \item \textbf{Land Cover:} Forests and shrublands are most fire-prone
    \item \textbf{Climate:} Mediterranean patterns strongly influence fire occurrence
    \item \textbf{Elevation:} Topography affects fire behavior through slope and aspect
    \item \textbf{Soil:} Subsurface factors affect vegetation and moisture
\end{itemize}

These insights inform preprocessing and feature engineering decisions in subsequent chapters.
